\begin{center}\begin{tikzpicture}[font=\small,>=stealth,line join=round]
\draw[gray!40,rounded corners=3pt] (0,0) rectangle (14.2,8.5);\node at (7.1,7.9) {易错四：同一个表达式 $y=x$，换了定义域就不同};
\node at (3.5,6.85){$D=[0,1]$：不是奇函数};\node at (10.6,6.85){$D=[-1,1]$：是奇函数};\begin{scope}[shift={(3.5,3.7)},x=1.5cm,y=1.5cm]\draw[->,gray] (-1.55,0)--(1.55,0) node[right]{$x$};\draw[->,gray] (0,-1.55)--(0,1.55) node[above]{$y$};\draw[blue,very thick] (0,0)--(1,1);\draw[red,dashed] (-1,-1)--(0,0);\fill[orange!85!black] (0,0) circle(2.2pt) node[below right=5pt] {$O$};\fill[orange!85!black] (1,1) circle(2.2pt) node[above right=5pt] {$(1,1)$};\draw[red,fill=white] (-1,-1) circle(2.5pt);\node[red,below left=5pt] at (-1,-1){$-1\notin D$};\end{scope}\begin{scope}[shift={(10.5,3.7)},x=1.5cm,y=1.5cm]\draw[->,gray] (-1.55,0)--(1.55,0) node[right]{$x$};\draw[->,gray] (0,-1.55)--(0,1.55) node[above]{$y$};\draw[blue,very thick] (-1,-1)--(1,1);\fill[orange!85!black] (1,1) circle(2.2pt) node[above right=5pt] {$(1,1)$};\fill[orange!85!black] (-1,-1) circle(2.2pt) node[below left=5pt] {$(-1,-1)$};\fill[orange!85!black] (0,0) circle(2.2pt);\end{scope}
\node[text width=12.8cm,align=center] at (7.1,.65) {两图同尺度；左图红虚线是缺失的对称像，不属于给定函数。};
\end{tikzpicture}\end{center}
